%Autore: rida
%Pēdējoreiz rediģēts 01.05.2022.
%Jautājumu un neskaidrību gadījumā rakstīt uz lienux13@gmail.com
%Sagatave paredzēta dažādu līmeņu darbu iesniegšanai Latvijas Universitātes Pedagoģijas, psiholoģijas un mākslas fakultātē, taču to droši vien viegli iespējams pielāgot citas fakultātes prasībām.
\documentclass[12pt]{article}
\usepackage{geometry}
\usepackage{gensymb} %grādu simbols
\usepackage[utf8]{inputenc}
\usepackage{tcolorbox}
\usepackage{multicol}
\usepackage{enumitem}
\usepackage{fancyhdr}
\usepackage{indentfirst}
\usepackage[bf,center]{titlesec}
\usepackage{tocloft}
\usepackage{hyperref}
\usepackage{tabularx}
\usepackage{ragged2e}
\usepackage{tocloft}
\usepackage{secdot}

\hypersetup{
	colorlinks=true,
	linkcolor=black,  
	urlcolor=blue
}

\urlstyle{rm}

\titlespacing*{\section}{0pt}{0pt}{1.5pt}
\titlespacing*{\subsection}{0pt}{0.5\baselineskip}{1.5pt}
\titlespacing*{\subsubsection}{0pt}{0.5\baselineskip}{1.5pt}
\titleformat{\section}{\large \bfseries \centering}{\thesection.}{1em}{}%section formatting
\titleformat{\subsection}{\normalsize \bfseries
	\centering}{\thesubsection.}{1em}{}%subsection formatting
\titleformat{\subsubsection}{\normalsize \bfseries
	\centering}{\thesubsubsection.}{1em}{}%subsubsection formatting
\geometry{a4paper, left=3cm, right=2cm, top=2cm, bottom=2cm}
\pagestyle{fancy}
\fancyhead{}
\fancyfoot{}
\fancyfoot[R]{\thepage}
\renewcommand{\headrulewidth}{0pt}
\linespread{1.3}
\setlength{\parindent}{1cm}
\setlength{\parskip}{0pt}
\renewcommand{\contentsname}{\hspace*{\fill}\bfseries\large Satura rādītājs\hspace*{\fill}} 
\renewcommand{\cfttoctitlefont}{\centering\bfseries\large}% Remove \bfseries from ToC title
\renewcommand{\cftsecfont}{}% Remove \bfseries from section titles in ToC
\renewcommand{\cftsecpagefont}{}% Remove \bfseries from section titles' page in ToC
\renewcommand{\cftsecleader}{\cftdotfill{\cftsecdotsep}}
\renewcommand\cftsecdotsep{\cftdot}
\renewcommand\cftsubsecdotsep{\cftdot}
\renewcommand{\cftsecaftersnum}{.}%dot after section number
\renewcommand{\cftsubsecaftersnum}{.}

\begin{document}
	
	\begin{titlepage}
		\thispagestyle{empty}
		\begin{center}
			\vspace*{2cm}
			\begin{large}\textbf{LATVIJAS UNIVERSITĀTE}
				
				\textbf{(FAKULTĀTES NOSAUKUMS) FAKULTĀTE}
				
				\vspace{1cm}
				\textbf{(NODAĻAS NOSAUKUMS) NODAĻA}
				
				\vfill
				\textbf{Darba nosaukums}
			\end{large}
			
			\vfill
			\textbf{Kursa nosaukums}
		\end{center}
		\vspace{4cm}
		\begin{flushright}
			Autore: \textbf{Tavs Vārds}
			
			Studenta apliecības nr.: tv22001

			%Darba vadītāja:
		\end{flushright}
		\vspace{1cm}
		\begin{center}
			RĪGA 2022
		\end{center}
	\end{titlepage}
	
	\tableofcontents
	
	\pagebreak
	\section{Virsraksts}
	Darba teksts
	
	Vēl darba teksts slkkjsbnfkffnksjbjs sfnkjf njkjf klmsdjnf enjnjfs knfdsnjfnjf nkfenjfjnfsnj jefwnjfnjf njfsnjfsnjfsnj nfdsje jnfksd;we klsfsop3w mklfkmfloew jwe pojew
	
	\subsection{Apakšnodaļa}
	Apakšnodaļas ievads
	
	\pagebreak
	\section*{Literatūras saraksts}
	\addcontentsline{toc}{section}{Literatūras saraksts}
	\begin{enumerate}
		\item Balson, M. (1996). \emph{Kā izprast klases uzvedību} (tulk. A. Zunde, M. Zelmenis, A. Puriņa). Lielvārde: Lielvārds. (Pirmoreiz publicēts 1992)
	\end{enumerate}
	
\end{document}